%!TEX program = xelatex
\documentclass[12pt, b5paper]{article}
\usepackage{ctex}

\usepackage{geometry}
\geometry{b5paper,left=2cm,right=1cm,top=2.5cm,bottom=2.5cm}
\usepackage{chemformula} %化学公式宏包
\usepackage[version=4]{mhchem} %化学公式宏包
\usepackage{siunitx} %数字，单位宏包
\sisetup{
	separate-uncertainty = true,
	inter-unit-product = \ensuremath{{}\cdot{}}
} %单位间加小圆点
\usepackage{tikz} %Tikz绘图包
\usetikzlibrary{cd}
\usetikzlibrary{chains}
\usetikzlibrary{arrows}
\usetikzlibrary{arrows.meta} %画箭头
\usetikzlibrary{commutative-diagrams} %交换图
\usetikzlibrary{positioning,automata}
\usepackage[all]{xy}
\usepackage{booktabs} %三线表
\usepackage{multirow} %合并单元格
\usepackage{makecell} %表格内强制换行
\usepackage{array}
\usepackage{booktabs} %控制表格行高
\usepackage{colortbl}  %彩色表格需要加载的宏包
\usepackage{amsmath}
\usepackage{unicode-math}
\usepackage{fontspec} %字体
\newcommand*{\cred}{\textcolor{red}}

\setmainfont{Times New Roman}
\setmathfont{XITS Math}

\setlength{\parindent}{0pt} %无缩进
\usepackage{setspace}
\usepackage{fancyhdr} %设置页码
\usepackage{lastpage} %获取总页数
\pagestyle{fancy} %fancyhdr宏包新增的页面风格
\renewcommand\headrulewidth{0pt} %取消页眉中的装饰分割线
\fancyhf{}
\cfoot{\footnotesize 第 \thepage\ 页(共 \pageref{LastPage}页)}%当前页 of 总页数

%微分符号
\makeatletter
\newcommand*{\dif}{\mathop{}\!\mathrm{d}}
\makeatother

%直立积分
\DeclareSymbolFont{EulerExtension}{U}{euex}{m}{n}
\DeclareMathSymbol{\euintop}{\mathop} {EulerExtension}{"52}
\DeclareMathSymbol{\euointop}{\mathop} {EulerExtension}{"48}
\let\intop\euintop
\let\ointop\euointop

\begin{document}
\begin{spacing}{1.625}

\centerline{{\Large 河北农业大学2022——2023学年第1学期}}

\centerline{\Large 参考答案及评分标准（B）卷}

\centerline{开课学院：\underline{\quad 理工学院 \quad} \quad 出\ \ 题\ \ 人：\underline{ \qquad \qquad 张斌 \qquad \qquad}}

\centerline{课\ \ 程\ \ 号：\underline{\quad H05981199 \ } \quad 课程名称：\underline{\quad 物理化学与胶体化学 \ \ }}

\bigskip


1. 简答题(15分)

要求至少列出3种生活中的热力学现象，并说明其原理。一些例子如下：

空调制冷原理；高海拔很难煮熟饭；内燃机；微波炉；霜升华成水蒸气等等。

\bigskip

\bigskip

2. 简答题(15分)

答题要点： 

(1)熵判据；吉布斯自由能判据；亥姆霍兹自由能判据。     $\hfill \cdots \cdots \cdots $(9分)

(2)用化学平衡判断反应进行的程度。   $\hfill \cdots \cdots \cdots $(6分)


\bigskip

\bigskip

3. 简答题(10分)

(1) C=S=2(\ch{CaCO3}和\ch{CO2})，P=2，又因\ch{CO2}压力一定，所以f'=C-P+1=1，这个自由度表示温度T可在一定范围内变化，而不出现CaO，即CaCO3不会分解。

$\hfill \cdots \cdots \cdots $(5分)

(2) 当\ch{CO2}压力确定不变，且\ch{CaCO3}(s)和CaO(s)共存时，C=S-R=3-1=2，P=3，则f '=C-P+1=0，即此时温度T也不能变化。也就是说，只有T、p都为某一定值不变时，\ch{CaCO3}(s)和CaO(s)才能共存。 $\hfill \cdots \cdots \cdots $(5分)

\bigskip

\bigskip

4. 主观题(20分)

\[
        \begin{tikzcd}[ampersand replacement=\&, column sep=small]
            \ce{H2O(g, 298K, 101.325kPa)} \ar[d, "\displaystyle{\Delta G_1}"']\ar[d, "\displaystyle{\text{定温可逆}}"] \ar[rrr, "\displaystyle{\Delta G}"] \& \& \&  \ce{H2O(l, 298K, 101.325kPa)} \\
            \ce{H2O(g, 298K, 3.168kPa)} \ar[rrr, "\displaystyle{\text{定温可逆}}"']\ar[rrr, "\displaystyle{\Delta G_2}"] \& \& \& \ce{H2O(l, 298K, 3.168kPa)} \ar[u, "\displaystyle{\text{定温可逆}}"']\ar[u, "\displaystyle{\Delta G_3}"]
              \end{tikzcd}
    \]

    
        $\Delta G_1 = nRT\ln \frac{p_2}{p_1} = -8586 \mathrm{J} \hfill  \qquad \cdots \cdots \cdots \cdots (4\text{分})$

        $\Delta G_2 = 0 \hfill \cdots \cdots \cdots \cdots (4\text{分})$

        $\Delta G_3 = \int_{p_1}^{p_2}V\dif p = \int_{p_1}^{p_2}V_m \dif p = nV_m(p_2 - p_1) = 1.8 \mathrm{J}  \hfill \cdots \cdots \cdots \cdots (4\text{分})$

        $\Delta G = \Delta G_1 + \Delta G_2 + \Delta G_3 = -8584J \hfill \cdots \cdots \cdots \cdots (4\text{分})$
    
    此过程$\Delta G < 0 $，所以是自发的。 $\hfill \cdots \cdots \cdots \cdots (4\text{分})$


\bigskip

\bigskip


5. 主观题(20分)

$\begin{array}{l}
\ln \frac{P_{2}}{P_{1}}=\frac{\Delta H_{\mathrm{m}}}{R}\left(\frac{1}{T_{1}}-\frac{1}{T_{2}}\right) \\
\ln \frac{31}{101.3}=\frac{41.05 \times 10^{3}}{8.314}\left(\frac{1}{373}-\frac{1}{T_{2}}\right) \\
T_{2}=342.5 \mathrm{K} \qquad \qquad \qquad \qquad \qquad \qquad \qquad \qquad \qquad \qquad \quad \cdots \cdots \cdots \cdots (6\text{分})\\
\ln \frac{k_{2}}{k_{1}}=\frac{-E_{a}}{R}\left(\frac{1}{T_{2}}-\frac{1}{T_{1}}\right) \\
\ln \frac{k_{2}}{k_{1}}=\frac{-85 \times 10^{3}}{8.314}\left(\frac{1}{342.5}-\frac{1}{373}\right) \\
\frac{k_{2}}{k_{1}}= 0.086  \quad \quad \qquad \qquad \qquad \qquad \qquad \qquad \qquad \qquad \qquad \qquad \cdots \cdots \cdots \cdots (6\text{分})\\
\text{一级反应}\quad \frac{k_{2}}{k_{1}}=\frac{t_{1}}{t_{2}} \\
\frac{10}{t_{2}} = 0.086 \\
t_{2}=116.28 min \\
\text{在珠穆朗玛峰顶上的沸水中煮熟一个鸡蛋需要116.28 min} \qquad \cdots \cdots \cdots \cdots (8\text{分})
\end{array}$

\bigskip

\bigskip

6. 主观题(20分)

评分标准：

使A的浓度远高于B的浓度，测定B的反应级数；使B的浓度远高于A的浓度，测定A的反应级数； $\hspace{15em} \cdots \cdots \cdots \cdots (10\text{分})$

测定A或B的反应级数可采用以下方法：

积分法；初速率法$\hspace{18em} \cdots \cdots \cdots \cdots (10\text{分})$

\end{spacing}
\end{document}
